% Skeleton of the first draft. Headings + [bracketed suggestions] + planned figures only.
% Suggestions are the lead's; prose is the author's. Numbers in brackets are provisional
% (no independent check yet) and come from reports/analyses/ (sources in paper/SOURCES.md).
\documentclass[11pt]{article}
\usepackage[margin=1in]{geometry}
\usepackage{amsmath,amssymb,graphicx,xcolor,hyperref,booktabs}
\usepackage[numbers]{natbib}

\newcommand{\sugg}[1]{\textcolor{blue!60!black}{[#1]}}
\newcommand{\figplan}[2]{%
  \begin{figure}[h]\centering
  \fbox{\parbox{0.9\linewidth}{\centering\small\textbf{Planned figure: #1}\\[2pt]\sugg{#2}}}
  \end{figure}}
\newcommand{\open}[1]{\textcolor{red!70!black}{[OPEN: #1]}}

\title{\sugg{Title options: ``Update Decomposition''; ``Reused or new? Testing whether fine-tunes reuse
pretrained mechanisms''; ``Scoring fine-tune diffing methods on fine-tunes of known construction''}}
\author{\sugg{authors}}
\date{}

\begin{document}
\maketitle

\begin{abstract}
\sugg{4--5 sentences, written last: question; definition; method + baselines on built pairs; headline 67M result
and its caveat (prior + steering); Qwen limit}
\end{abstract}

% ------------------------------------------------------------------
\section{Introduction}
\sugg{Motivation: post-training selects / re-weights pretrained mechanisms --- persona selection model, Jain et al.\
``thin wrapper'', alignment pretraining; evidence so far behavioural or activation-based}
\sugg{Why it matters: monitoring post-training; if reuse, interpretability of the base model transfers}
\sugg{Gap: no test of reuse \emph{inside} the model; no diffing method validated against a known reuse/new label}
\sugg{Approach in one line: define reuse; build fine-tunes whose answer is known; score methods against them}
\sugg{Contributions, bulleted: (1) definition(s) of reuse and where they break; (2) matched pairs of known
construction; (3) scorecard of re-gating vs cheap baselines on a 67M model; (4) limits at 1.7B}
\figplan{1 --- the question}{schematic: $W=\sum_c u_c v_c^\top+\Delta$, gate $\Gamma$, an update $\delta W$ and the
three places its effect can live (moved pieces / $\Delta$ / re-gated pieces); beside it one matched pair (copy vs
cipher) as the running example. Redraw of the log's overview diagram}

% ------------------------------------------------------------------
\section{Related work}
\subsection{Post-training as selection}
\sugg{persona selection; Jain et al.\ 2023/2024; Prakash et al.\ 2024 (entity tracking reused after fine-tune);
alignment pretraining; elicitation vs sharpening literature (lit\_review/elicitation\_vs\_sharpening.md)}
\subsection{Model diffing}
\sugg{fixed base-model SAE reports only features up/down (Wang et al.); crosscoders report ``new'' as an activation
direction, authors note they cannot tell new from old-in-new-places (Minder et al.); toy model: old feature in new
context reads as model-exclusive pair; no validation on fine-tunes of known status (paper\_notes.md)}
\sugg{do NOT say prior work could only report reuse --- false for crosscoders}
\subsection{Parameter decomposition}
\sugg{APD / SPD / VPD (Goodfire); rank-one subcomponents + causal-importance gate; weight-based diffing as an open
direction}
\subsection{Defining ``reuse''}
\sugg{causal abstraction (Geiger et al.); interchange interventions; Jain's possession + wrapper test; Prakash's
causal step}

% ------------------------------------------------------------------
\section{Setup}
\subsection{Models and decompositions}
\sugg{67M 4-layer Pile model + Goodfire's released VPD: 24 matrices, $\sim$39k subcomponents, $\sim$10k alive,
$\Gamma$ frozen; Qwen3-1.7B decomposition (new one only); one line each}
\subsection{Fine-tunes of known construction}
\sugg{the central design: matched pairs --- same inputs, push size, (prior); one side reuses something the model
half-does (copy, identity, case), other side is new by construction (cipher, random map, shuffled labels)}
\sugg{9 built pairs + 4 shuffled-label twins on 67M; hand-made weight edits (always-on, cue, table, XOR) as
known-mechanism checks; seeds per side}
\sugg{state once, plainly: ``reuse side'' = made to reuse; no circularity claim hidden}
\sugg{capacity checks: what the base model already does on each task}
\sugg{Table 1 instead of a figure: pairs $\times$ what is matched (inputs, size, prior, seeds)}
\open{Qwen pair library goes here or in \S7}

% ------------------------------------------------------------------
\section{What should ``reuse'' mean?}
\sugg{framed as a working definition (user, 11-17: assume one soon); the failed wordings become ``what a definition
must survive'', motivating the adopted one}
\open{which wording is adopted: placeholder until definition-lead lands one (retention, E196, may proxy update size)}
\subsection{Requirements}
\sugg{independent of any interpretability method; must look inside the model; total (always gives a reading); gives
reportable results; defensible}
\subsection{Candidate definitions}
\sugg{one line + one equation each; full wordings in App.~B}
\sugg{(i) Jain-style possession + wrapper; (ii) variable re-weighted / rewired / built (tentatively adopted 11-10,
cut-offs 0.8 / 0.5); (iii) soft intervention; (iv) simplest intervention, closure curves; (v) causal-rule retention}
\subsection{Results on the controls}
\sugg{Jain: no verdict on any task (possession fails everywhere)}
\sugg{(ii): ranks reuse above new in 5/5 pairs on 2 seeds, but absolute labels don't follow construction; ``built''
never read}
\sugg{(iii)/(iv): read hand-made edits correctly, give no verdict on real fine-tunes; closure curves separate 3/5
pairs on 4 seeds (copy/cipher, identity/random map, parity/fact), not case or repeat}
\sugg{(v): rebuilt circuit 0.97 vs 1; trained pairs: reuse side higher in 8/8 pair-seeds, but every fine-tune reads
$\geq 0.996$ (gaps $<0.003$) and the gaps follow the twins' MLP update size; MLP neurons only, misses attention (E196)}
\sugg{(v) held-model version (E197, E201): refit only output weights of the base model per site; built new head $S=0.13$, rebuilt circuit 0.56, every trained fine-tune $\geq 0.99$ except the fact side (0.65); reuse-minus-new gaps $<0.02$ on four pairs, 0.25 on twin/fact: separates built mechanisms, barely separates trained ones}
\sugg{common failure: real fine-tunes have no single named variable carrying the change; a search miss cannot be
told from a wording fault}
\figplan{2 --- definitions vs tests}{matrix: wordings (rows) $\times$ tests (reads hand-made edits / ranks matched
pairs / absolute labels follow construction / ever reads ``built''); cells pass/fail/no call, with the number}
\figplan{2b (optional) --- closure curves}{closure vs number of original variables used, reuse and new side of each
pair overlaid; from experiments/curve\_test}

% ------------------------------------------------------------------
\section{Measuring reuse on the 67M}
\subsection{The re-gating share}

The most effective measure of re-use we investigated is a VPD-based method we call the \textit{re-gating share}. The general idea is VPD's causal importance function $\Gamma$ should be better at knowing how to ablate existing mechanisms, than new ones. More precisely:
\begin{itemize}
	\item Let $W_0$ be the original weights and $W_1 = W_0 + \delta W$ the fine-tuned weights, with VPD subcomponents $C_0$.
	\item For some task relevant to the fine-tune, let $S_0, S_1$ be the scores before and after fine-tuning, with $G = S_1 - S_0$ being the gain in score on that task from the weight update. Let $a_0, a_1$ be the corresponding activations. Then let $S_{00}$ be the score of the components $C_0$ masked according to $g_0 = \Gamma(a_0)$ and $S_{01}$ the same but with masks $g_1 = \Gamma(a_1)$. 
	\item Then the \textit{re-gating share} is defined as $r := \frac{M_{01} - M_{00}}{G}$
\end{itemize}

Clearly, $0 \leq r \leq 1$, where $r$ intuitively measures how much of the gain from the fine-tune can be explained by $\Gamma$ turning on pre-existing components. In particular, $r = 0$ implies no benefit was gained by switching on components differently, while $r = 1$ implies all of the gain comes from re-gating.

Note that measuring $r$ requires a VPD decomposition of the base model, but does not require updating the components of re-training the decomposition in any way.

We studied the effectiveness of the regating share as a measure of elicitation by As shown in fig ??, the regating share performs remarkably well at distinguishing elicitation versus new mechanisms on the 67M
[what else can i say?]

\subsection{Baselines}

We compare the effectiveness of the re-gating share against a number of baselines for comparison. These range from relatively cheap methods such as SVD on the weight update, to more expensive techniques such as transcoders. More details can be found in Appendix \ref{sec:app_baseline}.

\subsection{Scorecard}
\sugg{share: reuse side higher on 9/9 built pairs, gaps 0.39--0.69; survives size matching (9/9; three pairs
naturally equal-size, gaps 0.52--0.62)}
\sugg{needs the trained decomposition: untrained SVD 2 same / 2 opposite / 5 no call}
\sugg{steering 4/3/2; neuron share $\le$3; update size 0 (new side always bigger)}
\sugg{agreement with \S4's definition on the pairs it reads}
\open{crosscoder row (E198): base-vs-itself control passed; planted new direction found by one latent (0.99 of the change) but relative norm 0.86, under the 0.9 ``fine-tune only'' bar; no real pair read; waits on the user (decision card P2.31, option a recommended)}
\open{share on hand-made edits (E200, run 1 of 2): share defined only on Table (0.26); not computable on always-on, cue, XOR in run 1; wait for run 2}
\figplan{3 --- scorecard}{(a) just re-gating (b) all methods (rows) $\times$ 13 pairs (cols), cell = reuse-minus-new gap coloured by
same/opposite/no call; twins in a separated block; redraw of method\_table/figures/gate\_scorecard\_1116.png}

% ------------------------------------------------------------------
\section{What the re-gating share measures}
\sugg{the main caveat: base-model prior + steering at strength 1 together reproduce all 9 direction calls}
\sugg{but gap sizes are not predicted by either (rank corr.\ 0.02, $-0.10$; 9 non-independent pairs)}
\sugg{one built pair where prior and steering disagree (E199): share follows steering (share gap 0.18, prior gap $-0.8$ decades); so the share is not just the prior, but nothing shown beyond steering; one pair, one construction}
\sugg{failure cases: shuffled-label twins (wrong way on 2/4); fine-tune to a different answer at the same trigger
reads like the real one; code-text fine-tune that hurt text still reads $\sim$0.8}
\sugg{is it a mechanism or the fine-tune leaking through $\Gamma$'s input? measured on VPD's masked stand-in, not the
fine-tuned model}
\sugg{granularity limit: a new circuit assembled from old pieces reads as re-gating}
\figplan{4 --- share vs familiarity}{per pair: share gap vs prior gap (and vs steering gap), size-matched points
overlaid, twins marked; shows direction shared, sizes not}

% ------------------------------------------------------------------
\section{Scaling to Qwen3-1.7B}
\open{filled last (user, 11-17); heading and claims wait on qwen-lead's results}
\sugg{share 0.99+ for real and random updates alike; log-odds version not computable}
\sugg{updates sit outside the subcomponents as much as random changes}
\sugg{frozen-neuron control decomposition: off the ceiling, 3/6 cells separate, 1/6 after size matching; built
switch-off reads mixed}
\sugg{RL vs SFT: one or two lines only --- RL raises GSM8K 0.78$\to$0.88, shuffled-reward null collapsed, no call
inside the decomposition; details App.~C}
\sugg{candidate reasons for saturation (decomposition too coarse at scale? MLP-only vs whole update?)}
\figplan{5 --- real vs random across scale}{share for real vs random updates on one axis: 67M, Qwen new
decomposition, Qwen frozen-neuron control}

% ------------------------------------------------------------------
\section{Discussion}
\sugg{what the evidence does and doesn't say about the re-weighting hypothesis (instruments tested, not the
hypothesis)}
\sugg{what a method that passes would need: positive control of known kind, size matching, beats prior + steering}
\sugg{recommendation: score diffing methods (incl.\ crosscoders) on built pairs}
\subsection{Limitations}
\sugg{one small model, one decomposition, one task family; $\le$9 non-independent pairs; labels by construction;
no significance tests; crosscoder comparison \open{pending E198}; cut-offs set on controls}
\subsection{Departures from the literature}
\sugg{per cited method: what we changed and why (Jain, Prakash, VPD, Goodfire neuron jobs)}

\section{Conclusion}
\sugg{3 bullets max}

\bibliographystyle{plainnat}
\bibliography{refs}

% ------------------------------------------------------------------
\appendix
\section{Carrying an update onto the decomposition}
\sugg{the original question; splits (rescaling, least-squares, adjoint) fail; teacher training keeps behaviour and
faithfulness on narrow fine-tunes, fails on a broad one; not tested on Qwen}
\figplan{A1}{behaviour kept vs adversarial faithfulness for do-nothing / splits / teacher training}
\section{Definition wordings in full}
\sugg{each wording, cut-offs, search, per-wording figure, what it read}
\section{Baseline implementations}\label{sec:app_baseline}
List of baselines
\section{RL and SFT on Qwen}
\sugg{GRPO recipe, null collapse, own-samples SFT, effective rank}
\section{Per-pair numbers and recipes}
\sugg{fine-tune recipe, seeds, tables behind every figure}

\end{document}
